\documentclass{article}
\usepackage[utf8]{inputenc}
\usepackage{amsmath}
\usepackage{amsfonts}
\usepackage{amssymb}
\usepackage{geometry}
\usepackage{tcolorbox}
\usepackage{hyperref}
\geometry{a4paper, margin=1in}

\title{Module 08: Production Pipelines \& Collaborative Git Workflows - Part 2}
\author{Cybernauts 2026 Datathon Campaign}
\date{May 2026}

\begin{document}

\maketitle

\begin{abstract}
During the high-stakes final rounds of elite datathons, pipelines must frequently ingest entirely hidden test datasets for live evaluation. A major engineering failure occurs when a feature pipeline breaks or leaks data because it was not built for true out-of-sample inference. This note details the mechanics of feature pipeline serialization using object persistence layers, ensuring absolute consistency across training and testing states.
\end{abstract}

\tableofcontents
\newpage

\section{The Live Inference Challenge}
In a standard offline environment, data scientists often manipulate training and testing sets simultaneously in memory. However, live onsite finals mirror production environments: your model is evaluated against a completely hidden test dataset that is processed entirely downstream.

\subsection{The Downstream Explosion}
If your feature engineering transformations rely on calculations derived from the training set, you cannot simply re-calculate them on the new test set. For example:
\begin{itemize}
    \item \textbf{Target Encoding:} Re-calculating target encodings on the test set is mathematically impossible because the test set target labels ($y_{\text{test}}$) are hidden.
    \item \textbf{Standard Scaling:} Re-computing $\mu$ and $\sigma$ on the test dataset shifts the coordinate system, meaning your model receives distorted features that do not match the patterns it learned during training.
    \item \textbf{High-Cardinality Categories:} Ingesting a test set containing categorical strings that never appeared in the training data will cause standard label encoders to crash with an \texttt{UnseenCategoryError}.
\end{itemize}

To survive these live evaluations, your preprocessing blocks must be decoupled from the data frames and treated as persistent, stateful software objects.

\section{Object Serialization (Pickling)}
Object serialization—commonly referred to as **Pickling** in Python—is the process of converting an in-memory object instance into a byte stream that can be saved directly to disk and reloaded later.

\subsection{State Preservation via \texttt{pickle} and \texttt{joblib}}
When a scikit-learn transformer calls \texttt{.fit()}, it computes structural parameters (e.g., the specific column medians for imputation, or the category-to-mean mappings for target encoding) and stores them as internal state variables ending in an underscore (like \texttt{scaler.mean\_}).

By saving these fitted object wrappers using serialization libraries like \texttt{pickle} or \texttt{joblib}, you lock those precise mathematical parameters into a file.

\section{The Operational Workflow Architecture}
To guarantee absolute consistency across data states and eliminate the risk of inference-time crashes, you must enforce a strict two-phase serialization protocol.

\begin{tcolorbox}[colback=blue!5!white,colframe=blue!70!black,title=The Pipeline Serialization Pattern]
\textbf{Phase 1: Training \& Export}
\begin{equation}
    X_{\text{train}} \xrightarrow{\texttt{transformer.fit\_transform()}} \text{State Extracted} \longrightarrow \text{Save to } \texttt{models/transformer.pkl}
\end{equation}
\textbf{Phase 2: Inference \& Transformation}
\begin{equation}
    \texttt{models/transformer.pkl} \longrightarrow \text{Load State} \xrightarrow{\texttt{transformer.transform()}} X_{\text{test}} \text{ (No Re-Fitting Allowed)}
\end{equation}
\end{tcolorbox}

\subsection{Phase 1: The Training Code Path}
During your offline cross-validation loop and final model training, fit your stateful transformers exclusively on the training split data. Immediately after fitting, serialize the object to your storage directory:
\begin{verbatim}
from sklearn.preprocessing import StandardScaler
import joblib

scaler = StandardScaler()
X_train_scaled = scaler.fit_transform(X_train) # Extracts and applies training mu and sigma
joblib.dump(scaler, "models/scaler.pkl")       # Serializes the fitted state to disk
\end{verbatim}

\subsection{Phase 2: The Inference Code Path}
When the platform provides the live hidden test set, your testing pipeline must load the serialized object from disk. You must **never** call \texttt{.fit()} or \texttt{.fit\_transform()} during this phase; instead, use the \texttt{.transform()} method to apply the frozen training statistics to the test data:
\begin{verbatim}
# Loaded inside your inference/submission script
fitted_scaler = joblib.load("models/scaler.pkl")
X_test_scaled = fitted_scaler.transform(X_test)  # Transforms test data using training parameters
\end{verbatim}
This workflow ensures that even if the test dataset contains extreme outliers or unexpected distribution shifts, the features are projected onto the exact same mathematical scale as your training data, preventing data leakage and stabilizing model performance.

\newpage
\section{Practice Problems}

\textbf{P1:} An engineer writes an onsite evaluation pipeline. During inference on the hidden test set, they write the following code block: \texttt{X\_test\_scaled = scaler.fit\_transform(X\_test)}. Explain why this script introduces a critical data leakage vulnerability.
\\
\textit{Answer: Calling \texttt{.fit\_transform()} recalculates the mean ($\mu_{\text{test}}$) and standard deviation ($\sigma_{\text{test}}$) directly from the hidden test distribution. This allows properties of the test set as a whole to influence the scaling of each individual row. In a production setting, this constitutes data leakage through distribution contamination, as the model's predictions for a given test sample will change depending on which other samples are included in the test set batch.}

\newline
\textbf{P2:} Write a robust Python function for a feature engineering script that handles missing value imputation. The function must include a boolean flag (\texttt{is\_training}). If \texttt{is\_training=True}, it must fit and save a scikit-learn \texttt{SimpleImputer} object to disk. If \texttt{is\_training=False}, it must load that exact file to transform the incoming data.
\\
\textit{Answer:}
\begin{verbatim}
import os
import joblib
import pandas as pd
from sklearn.impute import SimpleImputer

def process_imputation(df, features, model_dir="models", is_training=True):
    path = os.path.join(model_dir, "imputer.pkl")
    df_sub = df[features].copy()

    if is_training:
        os.makedirs(model_dir, exist_ok=True)
        imputer = SimpleImputer(strategy="median")
        df_imputed = imputer.fit_transform(df_sub)
        joblib.dump(imputer, path)
    else:
        if not os.path.exists(path):
            raise FileNotFoundError(f"Fitted imputer artifact missing at {path}!")
        imputer = joblib.load(path)
        df_imputed = imputer.transform(df_sub)

    return pd.DataFrame(df_imputed, columns=features, index=df.index)
\end{verbatim}

\newline
\textbf{P3:} Why is a standard python dictionary containing custom mapped string categories (e.g., \texttt{\{'low': 0, 'medium': 1, 'high': 2\}}) safe from standard pickling requirements compared to advanced target encoding routines?
\\
\textit{Answer: A custom ordinal mapping dictionary relies entirely on a static, pre-defined rule base that is independent of your data's target values or distribution. It does not calculate any statistical properties from your training data, so it can be safely imported as a static configuration metric inside a config script without risking data leakage or code crashes. Target encoders, by contrast, compute dynamic conditional probabilities from the training target arrays ($y_{\text{train}}$), making full object serialization mandatory.}

\end{document}
